\documentclass{article}\usepackage[margin=0.85in]{geometry}\begin{document}\section*{Same-root signed graph edge composition}
Six main declarations and twelve controls passed fresh Lean replay. Eighteen exact named standard-axiom readbacks passed. Nine modules returned zero through the Rust runner; the false equality returned native and runner exit one. Rust also performed inventory and hash acceptance. No Python executed.
\paragraph{Actual branch composition.} Two recursive find executions reach a common root. Their derived node equations reduce the original edge exactly as the native source. Opposite signs call the negated constraint. Equal signs use the literal conditional inverse of positive or negative two and then the anchor operation.
\paragraph{Global state and answer.} New anchors modify only the absent root entry. Existing anchors remain unchanged and generate the difference constraint. Both paths preserve Forest and every vertex/anchor equation; any returned scalar equals k. Inconsistent zero-coefficient assertions are explicitly represented and excluded under the true-equation premise.
\paragraph{Controls.} A four-vertex fixture retains reserved node zero and uses two nontrivial signed descendant finds. Both sign compositions, positive/negative denominator anchors, repeated-anchor recovery, consistent zero and inconsistent assertion outcomes were checked.
\paragraph{Premises and limits.} Field coefficients, nonzero two, initial Forest/Valid, a true incoming edge and returned signs in plus/minus one are explicit. Native inversion, word/sign encoding, counters, overflow, arrays/allocation, size metadata and compiled equivalence remain separate. No scientific execution occurred.
\paragraph{Custody.} Source-first plan, source hashes, native receipts and retained failures are included. Final sources await independent review after publication. Cached toolchain/dependencies remain trusted.
\end{document}
